Plancherel and \(|I|=1\) give the \(L^2(\mathbb R^n)\) to
\(L^2(I\times\mathbb R^n)\) norm bound \(\|\varphi\|_\infty\).
We next show an \(O(h^{-d/2})\) bound for \(p=\infty\).
Take a maximal \(c\sqrt h\)-separated net on \(S^{n-1}\) for
fixed small \(c>0\). Counting disjoint spherical balls of
half this radius bounds the number of centers by
\(O(h^{-d/2})\); the same volume comparison gives bounded
overlap of the caps of radius \(2c\sqrt h\).
A smooth partition on those caps is obtained from bumps in
\(|\omega-\omega_\nu|^2/h\), equal to one at distances at
most \(c\sqrt h\), by division by their sum. The denominator
is at least one. Multiply these cutoffs by \(\varphi(\eta)\).
In coordinates \(\eta=(\sqrt h\,z,\rho)\) with axis
\(\omega_\nu\), the resulting cutoff has fixed compact
support, with \(\rho\) bounded below. Its derivatives in
\((z,\rho)\) of any fixed order are uniformly bounded.
In fact each neighboring center has transverse coordinates
\(O(\sqrt h)\) and axial coordinate \(1+O(h)\) in this frame.
The transverse coordinates of \(\omega=\eta/|\eta|\)
are \(\sqrt h\,z/(\rho^2+h|z|^2)^{1/2}\), and
\((\omega\cdot\omega_\nu-1)/h\) is uniformly smooth on
this compact set after rationalizing the difference.
It follows that the arguments \(|\omega-\omega_\mu|^2/h\)
of the overlapping bumps have uniformly bounded derivatives.

Remove the linear phase \(T\rho/h\) from the multiplier phase.
The remaining phase factor is controlled by
\[
 \frac{T}{h}(\sqrt{\rho^2+h|z|^2}-\rho)
 =\frac{T|z|^2}{\sqrt{\rho^2+h|z|^2}+\rho}.
\]
It has uniformly bounded derivatives on the support for
\(T\in I\). Repeated integration by parts bounds the \(L^1\)
norm of the inverse Fourier transform in the rescaled
variables uniformly. Invertible linear frequency changes
preserve that kernel norm by the Jacobian, and multiplication
by the removed linear phase translates the kernel.
Thus every cap contributes \(O(1)\) to the spatial kernel
\(L^1\) norm. Their number proves the \(L^\infty\) endpoint
bound, uniformly in \(T\).

At \(p_c=2n/d\) use \ref{ms:thm-1-2}, with \(N=h^{-1}\) and
\(\log(2+N)\) comparable to \(\log(2/h)\).
Apply Riesz--Thorin to the linear operator
\(f\mapsto U(T)P_h f\) between the spatial and spacetime
Lebesgue spaces \cite{r-32}. For \(2\le p\le p_c\) the weight
of the critical estimate is
\[
 \frac{1/2-1/p}{1/2-1/p_c}=n(1-2/p),
\]
so the log exponent is \(d(1-2/p)\) with no power of \(h^{-1}\).
For \(p_c\le p\le\infty\) its weight is \(p_c/p\), giving
log exponent \(2/p\) and frequency exponent
\((d/2)(1-p_c/p)=d/2-n/p\). These are the stated exponents.
